\documentclass[12pt]{article}
\usepackage{amsmath, amssymb, amsthm}
\usepackage[margin=1in]{geometry}
\usepackage[skip=0.6em]{parskip}

\newtheorem{proposition}{Proposition}
\newtheorem{corollary}{Corollary}
\theoremstyle{remark}
\newtheorem{remark}{Remark}

\newcommand{\E}{\mathrm{E}}
\newcommand{\R}{\mathbf R}

\begin{document}

\begin{center}
{\Large Nice tests in the two-core model:\\[2pt] where an inadmissible example could live}\\[4pt]
{\small Notes on the question of whether all closed, convex, coordinatewise-lower\\
exact-size tests are admissible when $\mathcal C=\{1,2\}$ :: August 28, 2026}
\end{center}

\medskip

Throughout, $X\sim N_3(\mu,I_3)$, $\Lambda=\{\mu:\mu_3\le\max\{\mu_1,\mu_2\}\}$,
$\Lambda_0=\Lambda\cap(-\infty,0]^3$, and
$\Lambda_1=\{\mu\in\Lambda:\max\{\mu_1,\mu_2\}>0\}$. Call an acceptance region
$A$ \emph{nice} if it is closed, convex, coordinatewise lower, and
$P_0(A)=1-\alpha$. Call a direction $a\ge0$, $a\ne0$, \emph{feasible} if
$a_3\le\max\{a_1,a_2\}$ and \emph{bad} otherwise. Note that every nonzero
feasible direction belongs to $\Lambda_1$ (if $a_1=a_2=0$ then feasibility
forces $a_3\le0$, hence $a=0$). For a nice $A$, lowerness implies that the
size is attained at the origin, so any weak dominator $\psi$ of
$\phi_A:=I\{x\notin A\}$ satisfies, by the zero-mass argument of
Proposition~1 of the main note, $\E_0h=0$ with $h:=\psi-\phi_A$, and by the
first-moment argument there, $q:=\E_0[Xh(X)]=0$. Both arguments use only the
model, not the shape of $A$, so they are available for every nice test.

\section{Why examples are hard to find: three structural facts}

The difficulty is not bad luck. The two-core model destroys every standard
device for proving inadmissibility, for three related reasons.

\emph{First, the transfer cone is trivial.} Section~2 of the main note shows
that the only sample-space direction $d$ with $\mu'd\ge0$ for all
$\mu\in\Lambda_1$ is $d=0$. This is exactly the computation that makes
Stein--Birnbaum-type \emph{necessary} conditions for admissibility vacuous
here: those conditions require the rejection region to be monotone in
directions that raise every alternative likelihood ratio, and there are no
such directions. No test can be disqualified on local-monotonicity grounds.

\emph{Second, the convex hull of the alternative directions is all of
$\R^3$.} For instance $(-1,-1,0)=(1,-2,0)+(-2,1,0)$ with both summands in
$\Lambda_1$, and $(0,0,\pm1)$ arise similarly. This is why the first-moment
argument pins the entire vector $q=0$: a weak dominator of \emph{any}
exact-size test must match its rejection mass and all three first moments
under $P_0$. The usual dominators (transfers, gates) all fail this matching,
which is why they do not exist here.

\emph{Third, every nonnegative direction is either feasible or a negated
null direction.} If $a\ge0$ is bad, then $a_3>\max\{a_1,a_2\}\ge
\min\{a_1,a_2\}$, so $-a\in\Lambda_0$ (indeed
$-a_3\le-\min\{a_1,a_2\}=\max\{-a_1,-a_2\}$). In the Bayes-representation
heuristic for complete classes, a facet of an acceptance region can be
supported either by alternative prior mass along its normal or by null prior
mass along its negated normal. Since
$\overline\Lambda_1\cup(-\Lambda_0)\supseteq\R_+^3$, \emph{every} facet
orientation of a lower convex region is supportable from one side or the
other. There is no orientation-based obstruction to being a limit of Bayes
tests, hence no cheap route to inadmissibility.

The upshot: if an inadmissible nice test exists, the phenomenon must be
global and second-order, not directional. The next two results make this
precise by carving out large admissible classes.

\section{Feasibly supported regions are admissible}

Say $A$ is \emph{feasibly supported} if
\[
A=\bigcap_{a\in F}\{x:a'x\le h_A(a)\},
\qquad
F:=\{a\ \text{feasible}:h_A(a)<\infty\},
\]
that is, the supporting halfspaces with feasible normals already determine
$A$. For a polyhedron this holds if and only if every facet normal is
feasible.

\begin{proposition}[Feasible support suffices]\label{prop:feasible-support}
Every feasibly supported nice test is $\alpha$-admissible against all
level-$\alpha$ tests, including randomized competitors.
\end{proposition}

\begin{proof}
Let $\psi$ weakly dominate $\phi_A$ and $h=\psi-\phi_A$, so $\E_0h=0$,
$\{h>0\}\subseteq A$, $\{h<0\}\subseteq A^c$. Fix a nonzero $a\in F$; then
$a\in\Lambda_1$, so dominance along the ray $ta$, $t>0$, gives
$\E_0[e^{ta'X}h]\ge0$. Since $h^+$ is supported on
$A\subseteq\{a'x\le h_A(a)\}$, the tilt bound of the rectangle proof applies
verbatim: if $\E_0[h^-1\{a'X\ge h_A(a)+\varepsilon\}]>0$ for some
$\varepsilon>0$, then
$\E_0[e^{ta'X}h]\le e^{th_A(a)}\E_0h^+ -
e^{t(h_A(a)+\varepsilon)}\E_0[h^-1\{a'X\ge h_A(a)+\varepsilon\}]<0$ for
large $t$, a contradiction. Hence $h^-=0$ a.s.\ on $\{a'x>h_A(a)\}$. The
complement $A^c=\bigcup_{a\in F}\{a'x>h_A(a)\}$ is open, so by the
Lindel\"of property a countable subfamily covers it, and therefore $h^-=0$
a.s.\ on $A^c$. Since $\{h^->0\}\subseteq A^c$, this gives $h^-=0$ a.s.;
then $\E_0h=0$ forces $\E_0h^+=0$, so $h=0$ $P_0$-a.s., and mutual absolute
continuity of Gaussian shifts finishes the proof as in the main note.
\end{proof}

Note that the proof uses only the confinement step; the first-moment
identity is not needed. Two remarks. First, this is strictly weaker as a
hypothesis than the feasible-normal condition of the extension plan, which
demands that \emph{every} finite-support direction be feasible. Only the
directions needed to \emph{define} $A$ must be feasible; sums of facet
normals that leave the feasible set, such as $(1,1,2)$, are irrelevant.
Second, it settles a question the main note left open:

\begin{corollary}[The symmetric gate is admissible]
The symmetric gated region
$A_b^{\mathrm g}=\{x_1\le\kappa,\,x_2\le\kappa,\,x_1+x_3\le b,\,
x_2+x_3\le b\}$ has facet normals $e_1$, $e_2$, $(1,0,1)$, $(0,1,1)$, all
feasible, and is therefore admissible --- even though its finite-support
direction $(1,1,2)$ is infeasible and the extension-plan condition fails.
\end{corollary}

So the symmetric gate and the rectangle are \emph{both} admissible and
neither dominates the other, consistent with the power comparison in
Section~2 of the main note.

\section{One bad halfspace can be added for free}

\begin{proposition}[Feasible support plus one halfspace]
\label{prop:one-bad-facet}
Let $A=A_g\cap\{x:u'x\le c\}$ be nice, where $A_g$ is an intersection of
halfspaces with feasible normals and $u\ne0$ is arbitrary. Then $\phi_A$ is
$\alpha$-admissible.
\end{proposition}

\begin{proof}
As above, $\E_0h=0$ and now also $q=\E_0[Xh]=0$, so in particular
$\E_0[u'Xh]=0$. The confinement step of
Proposition~\ref{prop:feasible-support} gives $h^-=0$ a.s.\ on $A_g^c$, so
$h^-$ is supported on $A^c\cap A_g=\{u'x>c\}\cap A_g$, while $h^+$ is
supported on $A\subseteq\{u'x\le c\}$. With $M:=\E_0h^+=\E_0h^-$, if $M>0$
then $\E_0[u'Xh^+]\le cM$ and $\E_0[u'Xh^-]>cM$ (strictly, since $u'X>c$
a.s.\ on $\{h^->0\}$ and $P_0(u'X=c)=0$), whence $\E_0[u'Xh]<0$,
contradicting $u'q=0$. So $M=0$, $h=0$ $P_0$-a.s., and the conclusion
follows as before.
\end{proof}

This is the rectangle proof in its natural generality: the rectangle is
$\{x_1\le c_1,x_2\le c_2\}$ (feasible normals $e_1,e_2$) intersected with
the single bad halfspace $\{x_3\le c_3\}$. Two further corollaries are worth
recording.

\begin{corollary}
The redundant-only halfspace $\phi_3=I\{x_3>z_\alpha\}$ is admissible for
every $\alpha\in(0,1)$: here $A_g=\R^3$ and $u=e_3$, so
$q_3=0$ alone yields the contradiction. This removes the
Neyman--Pearson step from the proof of Proposition~2 of the main note, and
with it the $\alpha<1/2$ restriction flagged in my earlier comments.
\end{corollary}

\begin{corollary}\label{cor:one-cut}
The following exact-size regions are admissible: (i) any rectangle with one
additional \emph{feasible} cut, e.g.\
$\{x_1\le\kappa_1,\,x_2\le\kappa_2,\,x_3\le\kappa_3,\,a'x\le c\}$ with $a$
feasible (here $A_g$ collects the core faces and the cut, and the sole bad
halfspace is the $x_3$ face); and (ii) any region of the form
$\{x_1\le\kappa_1,\,x_2\le\kappa_2,\,a'x\le c\}$ with $a$ bad --- the two
core faces plus a single bad cut, with no $x_3$ face.
\end{corollary}

\begin{remark}[A rectangle plus a bad cut is \emph{not} covered]
\label{rem:not-covered}
Proposition~\ref{prop:one-bad-facet} does not apply to
$\{x_1\le\kappa_1,\,x_2\le\kappa_2,\,x_3\le\kappa_3,\,a'x\le c\}$ with $a$
bad and the cut binding: the $x_3$ face is itself a bad halfspace ($e_3$ is
infeasible), so this region has \emph{two} exposed bad facets and exceeds
the ``one arbitrary halfspace'' budget. Nor can the proof be repaired by a
cleverer choice of $u$: any separating functional $\ell$ must have
$\ell'x$ bounded below on the gap region, which contains points with
$x_3=\kappa_3+\varepsilon$ and $x_1,x_2\to-\infty$; this forces
$\ell_1=\ell_2=0$, i.e.\ $\ell\propto e_3$. But when the cut binds
($c<a_1\kappa_1+a_2\kappa_2+a_3\kappa_3$), the lobe $\{a'x>c\}\cap A_g$
contains points with $x_3<\kappa_3=\sup_Ax_3$, so $e_3$ fails as well. No
separating functional exists, and this region joins the frontier of
Section~4.
\end{remark}

\section{Where a counterexample must live}

Combining the two propositions: an inadmissible nice test must have at least
\emph{two} exposed halfspaces with bad normals, irreducibly --- the region
cannot be rewritten as feasible halfspaces plus a single extra one. (A
feasibly supported set cannot have a bad facet normal, since facet normals
of an intersection of halfspaces lie in the closed cone of the normals
used.) Moreover, the two bad facets must be arranged so that no single
linear functional separates $A$ from the gap region $A_g\setminus A$;
otherwise the proof of Proposition~\ref{prop:one-bad-facet} goes through
with that functional in place of $u$.

Two candidate families sit at this frontier. The simplest is the rectangle
with one bad cut of Remark~\ref{rem:not-covered},
\begin{equation}\label{eq:candidate-cut}
A^{\mathrm{cut}}
:=
\{x_1\le\kappa_1,\ x_2\le\kappa_2,\ x_3\le\kappa_3,\ a'x\le c\},
\qquad a\ \text{bad, cut binding},
\end{equation}
whose two bad facets are $e_3$ and $a$. The second is
\begin{equation}\label{eq:candidate}
A(\lambda)
:=
\{x_1\le\kappa,\ x_2\le\kappa,\ x_1+\lambda x_3\le b,\ x_2+\lambda x_3\le b\},
\qquad \lambda>1,
\end{equation}
with $(\kappa,b)$ calibrated to exact size; the gate normals
$(1,0,\lambda)$ and $(0,1,\lambda)$ are bad, both facets are exposed, and
the gap region has two lobes, one along each bad facet, which no single
linear functional can dominate simultaneously (the argument of
Remark~\ref{rem:not-covered} adapts). Every tool above fails on both
families, and I believe this is the true frontier: resolving admissibility
of, say, $A(2)$ or of $A^{\mathrm{cut}}$ with $a=(1,1,3)$ would either
produce the desired example or strongly indicate that none exists.

One observation tempers optimism about finding inadmissibility here. Both
families deform continuously into admissible tests. At $\lambda=1$,
$A(\lambda)$ is the symmetric gate --- admissible by
Proposition~\ref{prop:feasible-support}; as $\lambda\to\infty$ with $b$
rescaled, it converges to a rectangle --- admissible by
Proposition~\ref{prop:one-bad-facet}. Likewise $A^{\mathrm{cut}}$ tends to
the plain rectangle as the cut is relaxed ($c$ increased until it stops
binding) and to a region covered by
Corollary~\ref{cor:one-cut}(ii) as the $x_3$ face is relaxed
($\kappa_3\to\infty$). Inadmissibility, if present, would have to appear
and then disappear along these paths. Possible, but it would be a delicate
phenomenon.

\section{Conjecture and two ways to attack it}

\textbf{Conjecture.} In the two-core model, every closed, convex,
coordinatewise-lower exact-size test is admissible.

My reasons for leaning this way: the three structural facts of Section~1
(no disqualifying directions; all first moments pinned; every facet
orientation Bayes-supportable), plus the fact that every case resolved so
far --- all nonnegative halfspaces, all finite rectangles, $\phi_3$, the
symmetric gate, the regions of Corollary~\ref{cor:one-cut} --- has come out
admissible, plus the interpolation observation above. But the conjecture is
genuinely open on the families \eqref{eq:candidate-cut}
and~\eqref{eq:candidate}, and I would not bet heavily.

\textbf{Attack 1 (hunting: a linear-programming search).} Inadmissibility
of $\phi_A$ is the feasibility of a linear program in the perturbation
$h=\psi-\phi_A$: maximize $\sum_k \E_{\mu_k^*}h$ over a few probe
alternatives $\mu_k^*$, subject to $-\phi_A\le h\le1-\phi_A$,
$\E_\mu h\ge0$ for $\mu$ in a fine net of $\Lambda_1$ (include rays of
large norm --- these encode the tilt constraints that do the confining), and
$\E_\mu(\phi_A+h)\le\alpha$ for $\mu$ in a net of $\Lambda_0$. Discretize
$x$ on a truncated grid, impose the $1\leftrightarrow2$ symmetry to halve
the problem, and solve for $A(2)$. A robustly positive value as the nets
refine yields a concrete dominator candidate to verify analytically; a value
collapsing to zero is strong evidence for admissibility, and the dual
solution --- weights on the binding $\mu$'s --- reads off the prior sequence
for Attack~2. Either outcome is informative.

\textbf{Attack 2 (proving: a Bayes representation).} To prove admissibility
of $A(\lambda)$ directly, exhibit it as a suitable limit of Bayes tests and
apply a Blyth-type argument. The structure of Section~1 says which priors to
try: alternative prior mass escaping to infinity along the feasible facet
normals $e_1,e_2$ (these generate the core faces), and null prior mass
escaping along $-(1,0,\lambda)$ and $-(0,1,\lambda)$ --- both in $\Lambda_0$
since $\lambda>0$ --- to generate the bad facets from the null side, with
weights tuned so the limiting acceptance region is the polyhedron
$A(\lambda)$. The delicate point is that the null constraints far along
those rays are slack, so the weights must grow at the right exponential
rate; whether this can be done consistently for two bad facets at once is
exactly the same question the LP probes numerically.

\textbf{Literature worth revisiting.} The classical admissibility theory
for Gaussian means --- Birnbaum's convexity theorem, Stein's necessary
condition, Matthes--Truax for exponential families, and the
complete-class and one-sided-testing results of Marden, Perlman, and
Brown--Marden --- is where both the necessary-condition machinery and the
limit-of-Bayes sufficiency arguments live. As noted in Section~1, the
necessary conditions are vacuous in this model, which is itself the clearest
signal that the two-core problem sits in the regime where those papers'
sufficiency techniques, not their necessity techniques, are the relevant
ones. (Citations from memory; worth checking exact statements before
relying on them.)

\end{document}
